\hwpchapter{1}{데이터 이해 및 진단}{15}

\hwpsection{1.1 분석 대상 및 주요 변수}

KAMP 자원 최적화 AI 데이터셋(\texttt{okm\_augumented\_2021.csv})으로 생산조건에 따른 전력과 일 피크를 분석하였다. 한 공장의 시간별 생산량, 전력, 기상 기록이며 2021년 1월 1일부터 9월 14일까지 257일, 하루 24행씩 모두 6,168행과 18개 열로 이루어져 있다. 한 행에 그 시간의 15분 단위 전력 네 값(15분·30분·45분·60분 열)이 들어 있어 이를 펼치면 하루 96개, 모두 24,672개의 15분 슬롯이 된다. 예측 대상은 이 슬롯의 전력이며, 하루 96개 값의 최댓값을 일 피크(일 최대전력)로 정의하였다.

원자료를 직접 열어 확인한 내용은 다음과 같다. 결측은 풍속 3행, 강수량 1행, 공장인원 17행에만 있고 나머지 열에는 없다. 전력은 0--222 범위이며 0값이 74개 있다(1.3절). 생산량은 0--9,830이고 3,511행이 양수, 2,657행이 0이며, 257일 가운데 64일은 하루 종일 생산량이 0이다. \texttt{day}, \texttt{d}, \texttt{m} 열은 각각 요일(1=월요일, 7=일요일), 일, 월이며 6,168행 모두 날짜 열에서 얻는 값과 같다. 열별 형태와 이 보고서에서의 사용 방식은 표~\ref{tab:ch1-vars}에 정리하였다.

{\hwptablefont
\setlength{\tabcolsep}{4pt}
\begin{longtable}{L{2.2cm}L{3.3cm}L{4.0cm}L{5.7cm}}
\caption{원자료의 열과 사용 방식.}\label{tab:ch1-vars}\\
\toprule
원자료 열 & 형태·범위 & 공정상 의미 & 사용 \\
\midrule
\endfirsthead
\caption[]{원자료의 열과 사용 방식(계속).}\\
\toprule
원자료 열 & 형태·범위 & 공정상 의미 & 사용 \\
\midrule
\endhead
\midrule
\multicolumn{4}{r}{\footnotesize 다음 쪽에 계속}\\
\endfoot
\bottomrule
\endlastfoot
날짜 & 정수(YYYYMMDD). 2021-01-01--09-14, 257일 & 관측일. 하루 24행(시간별) & 예측일 식별. 시간순 분할과 요일·월·공휴일 계산의 기준 \\
시간 & 정수 0--23(48행은 이 범위 밖의 값) & 하루 안의 시각(시) & 슬롯 시각 부여. 7월 13일과 15일의 48행은 값이 어긋나 행 순서로 대체 \\
15분·30분·45분·60분 & 정수 0--222(kW로 해석). 0값 74개 & 공장 전체의 15분 평균전력. 한 행에 그 시간의 네 슬롯 & 평가 목표(하루 96슬롯, 일 피크). C-BAT 입력: 과거 전력(학습 목표, 기준일 곡선). 비교모형 입력: 과거 전력(M2·LightGBM(튜닝)은 시차·요약 변수, B0·B0\_kind는 기준일 곡선) \\
평균 & 정수 0--208 & 네 전력값 평균의 반올림(6,168행 모두 일치) & 제외: 전력에서 계산된 값 \\
생산량 & 정수 0--9,830. 양수 3,511행, 0 2,657행 & 그 시간의 생산 실적. 공정이 돌았는지와 생산 규모 & 실적을 생산계획으로 간주. C-BAT 입력: 시간별 24개 값에서 생산 여부·정규화 로그 생산량·가동유형 산출. 비교모형 입력: M2·LightGBM(튜닝)은 시간별 값·일 합계·생산시간, B0\_kind는 가동유형으로 기준일 선택 \\
기온·풍속·습도·강수량 & 기온 $-12.0$--33.4, 풍속 0.0--7.6, 습도 8--98(정수), 강수량 0.0--122.4(누적). 풍속 3행, 강수량 1행 결측 & 공장 밖의 기상 & C-BAT 제외: 공장 위치를 알 수 없어 기상 예보를 연결할 수 없다. 비교모형 입력: M2·LightGBM(튜닝)은 전일 기온·습도·풍속 평균과 전일 강수 합계만 입력. B0·B0\_kind는 미사용 \\
전기요금(계절) & 실수 세 값: 109.8(1--2월), 167.2(3--5월, 9월), 191.6(6--8월) & 계절별 요금 단가. 한 달 안에서는 일정 & 제외: 예측 입력이 아니다. 비용 분석은 시간대를 구분하는 한국전력 산업용(을) 2021년 요금표로 계산 \\
\texttt{day}·\texttt{d}·\texttt{m} & 정수. day 1--7, d 1--31, m 1--9 & 요일, 일, 월. 하루 24행이 같은 값 & 제외: 날짜 열과 중복. 요일·월은 날짜에서 다시 계산하여 사용 \\
공장인원 & 실수. 17행 결측(전력 네 값이 모두 0인 17개 시간) & 생산량을 한 시간 안 전력 네 값의 합으로 나눈 값. 실제 작업자 수가 아님 & 제외: 목표 전력이 계산에 들어 있어 정보 누수 \\
인건비 & 1.0(2,295행) 또는 1.5(3,873행) & 시각별 인건비 할증 배수. 09--17시 2,313행 중 2,295행이 1.0, 그 밖의 시각은 모두 1.5 & 제외: 예측 입력이 아니며 분석에도 쓰지 않았다 \\
\end{longtable}}
\hwpnote{C-BAT는 이 보고서가 제안하는 모형으로, 시간별 생산계획을 바탕으로 한 조건부 베이지안 전달모형이며 2.1절에서 설명한다. 전력은 kW로 해석하였다. 자료 설명에 단위가 없는 생산량과 기상 열은 값 그대로 사용하였다. 비교모형은 B0(7일 전 전력 곡선), B0\_kind(같은 가동유형의 기준일 곡선), M2(LightGBM 기본), LightGBM(튜닝), 기존 BAT이며 구성은 2장에서 설명한다. 기존 BAT는 C-BAT와 같은 변수를 쓴다.}

원자료에 없는 파생 변수는 표~\ref{tab:ch1-derived}에 정리하였다. 가동유형은 하루에 생산량이 양수인 시간 수로 정하므로 생산계획에서 바로 계산되고, 예측 시점에 알 수 있다.

\begin{table}[htbp]
\hwptablefont
\setlength{\tabcolsep}{4pt}
\caption{원자료에서 만든 파생 변수.}\label{tab:ch1-derived}
\begin{tabular}{L{2.5cm}L{6.4cm}L{6.3cm}}
\toprule
파생 변수 & 정의 & 사용 \\
\midrule
생산 여부 & 시간별 생산량이 양수이면 1, 아니면 0 & C-BAT 입력(생산 여부 채널). M2·LightGBM(튜닝) 입력(슬롯별 생산 여부) \\
정규화 로그 생산량 & $\log(1+q)/\log(1+q_{\max})$. $q$는 시간별 생산량, $q_{\max}$는 학습 기간의 시간별 최대 생산량 & C-BAT 입력(로그 생산량 채널). M2·LightGBM(튜닝)은 정규화하지 않은 생산량 원값 사용 \\
하루 생산시간·일 총생산량 & 생산량이 양수인 시간 수, 시간별 생산량의 합 & M2·LightGBM(튜닝)의 일 피크 모형 입력 \\
가동 여부 & 그날 생산량의 합이 0보다 크면 가동(시각 오류일 2일은 가동으로 표시) & C-BAT 기준일 선택. M2·LightGBM(튜닝) 입력. B0는 7일 전 곡선이 없을 때 대체 기준일 선택에만 사용 \\
가동유형 & 0(0시간)·1(1--12시간)·2(13--19시간)·3(20시간 이상) & C-BAT 입력(유형별 기본 곡선, 가동 상태별 잡음). B0\_kind 기준일 선택. M2·LightGBM(튜닝)은 가동유형 대신 가동 여부와 생산시간 사용 \\
날짜유형 & 평일(월--금), 토요일, 일요일 & C-BAT 기준일 선택. M2·LightGBM(튜닝) 입력. B0\_kind 기준일 선택 \\
요일·월 & 날짜에서 계산(계절은 월에서 계산) & M2·LightGBM(튜닝) 입력. C-BAT 미사용 \\
공휴일 & 2021년 8일(1월 1일, 2월 11--13일, 3월 1일, 5월 5일, 5월 19일, 8월 16일)을 별도 목록으로 관리 & M2·LightGBM(튜닝) 입력. 비용 분석(요금 시간대 판정). C-BAT 미사용 \\
\bottomrule
\end{tabular}
\end{table}
\hwpnote{기존 BAT는 이 표의 C-BAT 입력과 같은 파생 변수(생산 여부, 정규화 로그 생산량, 가동유형, 가동 여부, 날짜유형)를 쓴다.}

정리하면 C-BAT는 예측일의 시간별 생산계획, 예측일 이전의 전력, 가동 여부와 날짜유형만 입력하며 기상, 공휴일, 요일, 월은 쓰지 않았다. M2와 LightGBM(튜닝)은 여기에 달력 정보(요일·월·공휴일)와 전일 기상·전일 생산 요약값을 더하였다.

\hwpsection{1.2 제조공정 상태의 해석}

본 자료는 공장 전체의 생산량과 전력을 집계한 자료이다. 생산시간과 물량 집중에 따른 부하 변화를 분석하는 데에는 적합하나, 개별 설비의 가동 여부, 제품 종류, 교대조 정보가 없어 특정 설비나 공정이 피크를 일으킨 원인은 직접 확정하지 않았다. 하루 종일 생산량이 0인 날에도 전력은 관측되어(비가동일의 일 피크 중앙값 27 kW), 생산 실적이 없는 상태와 설비 전체의 정지를 구분하였다.

생산량이 양수인 시간 수에 따라 하루를 가동유형 0(비가동, 0시간)·1(1--12시간)·2(13--19시간)·3(20시간 이상)으로 구분하였다. 복사일, 시각 오류일, 결측이 있는 날을 제외한 1--8월 완전관측일 117일에서 가동유형별 일 피크의 중앙값은 각각 27, 112, 186, 193 kW였다(표~\ref{tab:ch1-kind}). 생산시간이 긴 유형에서 일 피크가 높았으나 증가 폭은 일정하지 않았고(가동유형 2와 3의 차이는 7 kW), 유형별 범위가 서로 겹쳐 가동유형만으로 일 피크가 정해지지는 않았다. 이에 따라 가동유형으로 운영 상태를 구분하고 시간별 생산량과 과거 전력 곡선을 함께 쓰도록 모형을 구성하였다(2장).

\begin{table}[htbp]
\centering
\hwptablefont
\caption{가동유형별 일 피크(1--8월 완전관측일 117일, kW).}\label{tab:ch1-kind}
\begin{tabular}{L{4.2cm}rrr}
\toprule
가동유형 & 일수 & 중앙값 & 범위 \\
\midrule
0 (비가동, 0시간) & 24 & 27 & 25--116 \\
1 (1--12시간) & 18 & 112 & 104--197 \\
2 (13--19시간) & 22 & 186 & 172--222 \\
3 (20시간 이상) & 53 & 193 & 119--222 \\
\bottomrule
\end{tabular}
\end{table}
\hwpnote{비가동일의 최댓값 116 kW는 1월 2일 하루이며 나머지 23일은 25--31 kW였다. 생산시간은 생산 실적이 기록된 시간 수이며 설비 운전시간이나 교대조 수를 뜻하지 않는다. 유형별 차이는 관측된 연관성으로 해석하였다.}

개발 구간(2021년 7월 7일부터 8월 31일까지를 14일씩 나눈 검증 구간 f1--f4, 1.4절)에서 일 피크를 평가하는 51일 가운데 고피크일은 13일이다. 고피크일은 일 피크가 해당 학습 기간 가동일 일 피크의 90\% 분위수(C90)를 넘은 날이며, C90은 검증 구간별로 198.0--211.0 kW였다. 고피크일 13일은 모두 13시간 이상 가동한 날(가동유형 2가 3일, 3이 10일)이었고, 모두 일 피크가 나온 슬롯에 생산이 있었다. 개발 구간 가동일 37일의 첫 최댓값 시각은 08시가 10일, 11시가 9일로 두 시각에 절반가량(19일) 몰렸으며, 고피크일 13일 중 10일(08시 6일, 11시 4일)이 이 두 시각에 피크를 기록하였다(그림~\ref{fig:ch1-peakhours}). 피크가 생산이 있는 특정 시각에 생기므로 생산계획을 입력으로 받아 시각별 부하를 예측하는 모형이 필요하다. 이 관계는 사후 연관이며 표본이 작다(고피크일 13일).

\begin{figure}[htbp]
\centering
\includegraphics[width=\linewidth]{peak_hours.pdf}
\caption{개발 구간 가동일 37일의 첫 피크 시각 분포. 짙은 색은 고피크일 13일, 옅은 색은 그 밖의 가동일 24일이다(그림 속 범례는 영어: High peak, Other op.).}\label{fig:ch1-peakhours}
\end{figure}

\hwpsection{1.3 데이터 진단 및 전처리}

전처리에서는 기록 오류, 중복, 정보 누수를 통제하였다. 높은 전력값은 피크 분석에 필요한 관측값이므로 크기만을 기준으로 제거하거나 상한을 두지 않았다.
\begin{hwplist}
\item 시계열 정렬: 시간별 행의 전력 4개 열을 15분 슬롯으로 펼쳐 단일 시계열로 만들고 각 슬롯에 시작 시각을 부여하였다. 시간별 생산량과 기상은 그 시간의 네 슬롯에 같은 값으로 연결하되, 일 총생산량은 원래 시간별 값을 합산하여 중복 집계를 막았다.
\item 시각 오류: 7월 13일과 15일의 48개 행에서 기록된 시각 열이 행 순서와 맞지 않았다. 두 날의 전력과 생산량을 결측 처리하고 오류 표지를 유지하여 학습과 평가에서 제외하였다. 자료 정렬에는 행 순서를 썼으며, 이를 실제 시각의 복원으로 보지는 않았다.
\item 전력 결측: 전력 0값 74개(8월 28일 26개, 8월 29일 46개, 9월 8일 2개)를 계측 손실로 보고 결측 처리하였다. 피크 왜곡을 막기 위해 목표 전력은 보간하지 않았다. 슬롯 오차는 관측된 슬롯만으로 계산하고, 일 피크 평가는 결측이 4슬롯 이하인 날로 제한하였다.
\item 중복 제거: 하루 96개 전력값이 소수점까지 같은 날을 한 묶음으로 보고 한 날만 원본으로 남겨 115일을 완전 복사일로 표시하였다. 원본은 시간순 첫 날로 하되, 첫 날이 생산 없이 40 kW를 넘는 슬롯을 48개 이상 가지면 생산이 있는 날을 원본으로 삼았다. 같은 복사쌍의 기온과 습도 시계열은 115쌍 모두 서로 달라, 전력 곡선만 복사해 넣은 증강 자료로 판단하였다. 나머지 날 가운데 40 kW를 넘는 슬롯이 다른 날의 같은 슬롯과 20개 이상 일치한 7일은 편집 복사일로 분류하였다. 복사일 122일(257일의 47\%)의 전력은 결측 처리하여 학습, 평가, 과거 전력 참조에서 제외하였다.
\item 정보 누수 변수: 공장인원은 생산량을 한 시간 안 전력 네 값의 합으로 나눈 값과 일치하였고(유효 6,151행, 최대 절대오차 $5\times10^{-9}$), 평균은 전력 네 값의 평균을 반올림한 값과 6,168행 모두 일치하였다. 예측 대상이 계산에 들어 있는 두 열을 입력에서 제외하였으며, 공장인원을 실제 작업자 수로 해석하지 않았다.
\item 기상정보: 풍속 결측 3개 시간은 선형보간하고, 누적 강수량은 날짜 경계를 포함한 연속 차분으로 증분화하였다. 음의 차분은 누적값이 초기화된 것으로 보아 현재값을 사용하였다. 강수량 결측 1개 시간은 증분을 0으로 두고 직전 유효 누적값을 유지하였다.
\item 파생변수: 생산량에서 생산 여부, 하루 생산시간, 일 총생산량, 가동유형을 만들었다. C-BAT에는 생산 여부와 정규화한 로그 생산량을 구분하여 입력하였다. 요일, 월, 평일·토·일은 날짜에서 만들고 공휴일은 별도 목록으로 관리하였다. 정규화 기준과 전력 기본 수준은 학습 기간에서만 산출하였다.
\end{hwplist}
\hwpnote{전력 0값의 원인이 실제 정전인지 계측 누락인지는 확인되지 않았다. 결측이 적은 날에도 관측되지 않은 슬롯에서 일 피크가 발생했을 가능성은 남아 있다.}

\hwpsection{1.4 예측 시점의 정보 사용 범위와 검증 설계}

매일 00:00 시점에 당일 96슬롯의 전력을 예측하도록 설정하였다. 입력은 당일의 시간별 생산계획, 사전에 알 수 있는 달력 정보, 예측일 이전에 관측된 전력으로 한정하였다. 모든 모형은 예측일과 그 이후의 전력을 가린 자료로 예측하므로 당일과 미래의 전력값은 입력에 들어가지 않는다.

제공된 자료에는 생산 실적만 있고 사전 생산계획은 없다. 생산 실적을 전날 확정된 생산계획으로 간주하였으므로, 이 보고서의 모든 성능은 시간별 생산계획을 미리 안다는 가정에서 나온 값이다. 현장 적용 시에는 실제 계획 자료를 연결하고 계획과 실적의 차이가 미치는 영향을 확인해야 하며, 이 영향은 3장의 입력 오차 실험에서 살핀다.

복사일이 섞인 자료에서 날짜를 무작위로 나누면 시험일의 사본이 학습에 들어간다. 이를 확인하기 위해 9월 이전에 관측된 241일을 날짜 단위 무작위 5-fold와 LOEO(leave-one-event-out)로 비교하였다. LOEO는 같은 전력 곡선을 공유하는 날의 묶음(복사 사건)을 함께 제외하며, 사건은 126개(복사군 45개, 단독일 81개)이다. M2의 MAE는 무작위 분할에서 14.91 kW, LOEO에서 16.64 kW로 차이는 $-1.73$ kW이고 95\% 신뢰구간은 0을 포함하지 않는다(표~\ref{tab:ch1-leak}). 기존 BAT는 두 분할에서 거의 같았다. 무작위 분할은 M2를 약 10\% 낙관적으로 보이게 하는 누수를 만들었다.

\begin{table}[htbp]
\centering
\hwptablefont
\caption{복사일을 포함한 자료의 누수 진단(MAE, kW, 관측 241일). 차이는 무작위 5-fold에서 LOEO를 뺀 값이다.}\label{tab:ch1-leak}
\begin{tabular}{lrrr}
\toprule
모형 & 무작위 5-fold & LOEO & 차이 [95\% 신뢰구간] \\
\midrule
M2 (LightGBM 기본) & 14.91 & 16.64 & $-1.73$ [$-3.63$, $-0.04$] \\
기존 BAT & 16.37 & 16.32 & $+0.05$ [$-0.22$, $+0.42$] \\
\bottomrule
\end{tabular}
\end{table}
\hwpnote{신뢰구간은 날짜 단위 짝지은 부트스트랩(2,000회)으로 구하였다. 진단 실행은 시간을 줄이려 기존 BAT를 1,000회 표집(초기 500회 제외), M2를 100회 부스팅으로 줄인 설정이며, 누수의 정도를 보는 값이지 전방 예측 성능의 추정치가 아니다.}

이 결과에 따라 모든 모형 비교는 복사일을 목표에서 뺀 시간순 분할로 하였다. 개발 구간의 검증 구간 f1--f4는 2021년 7월 7일부터 8월 31일까지 14일씩 네 개(f1 7월 7--20일, f2 7월 21일--8월 3일, f3 8월 4--17일, f4 8월 18--31일)이며, 각 검증 구간의 모형은 그 이전 자료만으로 학습한다. 검증 첫날의 전날은 간격일로 두어 목표에서 빼고 첫 검증일의 입력으로만 쓴다. 검증 구간과 같은 복사 사건에 속한 학습일은 제거(purge)하는 규칙을 두었으나 실제로 해당하는 날은 없었다. 학습 기간에는 복사일이 많아(표~\ref{tab:ch1-copies}), f1의 학습 기간은 달력상 186일이지만 복사일 121일(65\%)을 빼면 65일이다.

\begin{table}[htbp]
\centering
\hwptablefont
\caption{검증 구간별 학습 기간의 달력일과 복사일(일).}\label{tab:ch1-copies}
\begin{tabular}{lrrr}
\toprule
검증 구간 & 학습 달력일 & 복사일 & 복사 제외 \\
\midrule
f1 & 186 & 121 & 65 \\
f2 & 200 & 121 & 79 \\
f3 & 214 & 122 & 92 \\
f4 & 228 & 122 & 106 \\
9월 시험 & 242 & 122 & 120 \\
\bottomrule
\end{tabular}
\end{table}

개발 구간 56일 가운데 시각 오류일 2일(7월 13일, 15일)과 복사일 1일(7월 30일)은 목표가 없어 채점에서 빠지므로, 유효하게 채점되는 날은 53일(5,016슬롯)이다. 이 중 8월 28일과 29일은 결측이 5슬롯 이상이어서 일 피크 평가일은 51일이다.

9월 1--14일은 최종 평가를 위한 9월 시험 구간으로 봉인하였다. 기본 로더가 이 구간의 목표와 입력을 모두 가리며, 이전 개발 단계에서 한 번 열린 적이 있어 완전한 미사용 구간은 아니다. C-BAT는 개발 구간에서 확정하여 동결한 뒤 9월에 적용하였고, 9월 결과로 구조나 하이퍼파라미터를 고르지 않았다. 주 성능은 개발 구간과 9월 시험 구간을 합친 67일·6,358슬롯(피크 평가일 65일)의 시간순 평가(합산 평가)로 제시하며, 합산 평가의 정의는 2.2절에서 다룬다. 모델 선택, 절제, 오류분석은 개발 구간에서, 9월 단독 결과는 보조로 사용하였다.
